\section{Identificando o problema}
O ensino de Química, especialmente no que se refere ao domínio da Tabela Periódica, enfrenta uma persistente dificuldade de aprendizagem associada à predominância de práticas pedagógicas tradicionais baseadas na memorização mecânica, em detrimento de abordagens que promovam o engajamento ativo dos estudantes [1], [2]. Nesse contexto, a periodicidade química é frequentemente percebida como um conteúdo abstrato e de alta complexidade, sendo abordada por meio de estratégias de “decoreba” que não favorecem a construção de uma compreensão significativa de suas relações estruturais e funcionais [2].
Evidências empíricas demonstram a gravidade dessa lacuna: estudos diagnósticos revelam que estudantes de nível médio frequentemente falham em atingir o limiar básico de proficiência de 75% em testes de compreensão da Tabela Periódica [2]. No cenário brasileiro, essa dificuldade é corroborada pelos resultados do PISA 2022, que indicam que apenas 45% dos alunos atingem o nível mínimo de proficiência em ciências [3]. 
Além das dificuldades cognitivas, a predominância de métodos expositivos e centrados no professor impacta negativamente a motivação discente e a retenção de conhecimento a longo prazo [1]. A literatura aponta que a aprendizagem significativa está diretamente relacionada à satisfação de necessidades psicológicas básicas, como autonomia, competência e pertencimento [2]. No entanto, o modelo tradicional de ensino tende a limitar a participação ativa do estudante, reduzindo sua agência no processo de aprendizagem e contribuindo para o desinteresse e a baixa retenção de conteúdos [4]. 
